\section{Preliminaries}
\begin{definition}[KKT Point and $(\varepsilon_1, \varepsilon_2)$-KKT Point]
\label{def:kkt}
A point $\bx^* \in \calX$ is called a \emph{KKT point} 
of problem \eqref{eq:nonlinear-composite} if
\begin{equation}\label{eq:kkt}
    0 \in \nabla f(\bx^*) + \partial g(\bx^*) 
    + J_{\bc}(\bx^*)^\top \partial h\big(\bc(\bx^*)\big)
    + \mathcal{N}_{\calX}(\bx^*).
\end{equation}
Moreover, for $\varepsilon_1, \varepsilon_2 > 0$, a point 
$\bx^* \in \calX$ is called an 
\emph{$(\varepsilon_1, \varepsilon_2)$-KKT point} if there exists 
$\by^* \in \mathbb{R}^m$ such that:
\begin{align}
    \mathrm{dist}\Big(- \mathcal{N}_{\calX}(\bx^*),\; \nabla f(\bx^*) + \partial g(\bx^*) 
    + J_{\bc}(\bx^*)^\top \partial h\big(\by^*\big)\Big) 
    &\leq \varepsilon_1, \label{eq:eps-kkt-stationarity}\\
    \|\bc(\bx^*) - \by^*\| 
    &\leq \varepsilon_2. \label{eq:eps-kkt-feasibility}
\end{align}
\end{definition}
\begin{remark}
Condition \eqref{eq:kkt} is derived via the chain rule under 
Assumption~\ref{ass:CQ}. In the block-separable formulation 
\eqref{eq:block-separable}, introducing a dual multiplier 
$\lambda^* \in \mathbb{R}^m$, the KKT condition 
\eqref{eq:kkt} is equivalently expressed as:
\begin{align}
    0 &\in \nabla f(x^*) + \partial g(x^*) 
        + J_{\bc}(x^*)^T \lambda^*, \\
    0 &\in \partial h(y^*) - \lambda^*, \\
    0 &= \bc(x^*) - y^*,
\end{align}
which reduces to \eqref{eq:kkt} upon eliminating $\lambda^*$.
\end{remark}

\subsection{Main Assumptions}
To establish the ensuing analysis, we introduce the following standard assumptions for problem

\begin{assumption}[Constraint Qualification]\label{ass:CQ}
The following constraint qualification holds for problem 
\eqref{eq:nonlinear-composite}:
\[
    \text{[CQ]} \qquad 
    \nabla \bc(\bx)^T y = 0, \quad 
    y \in \partial^\infty h\big(\bc(\bx)\big) 
    \implies y = 0.
\]
\end{assumption}
